\section{API 调用与模型行为数据采集}

\subsection{实验样本范围}

本轮 API 实验不使用全部候选池，而是选取 5 组代表性 Prompt 来源：\path{data/prompts.csv} 作为 baseline，加上当前重点策略 \path{real_world_archetype_matrix}、\path{halton_space_filling}、\path{dimensionless_ratio_design} 和 \path{constraint_boundary_solver}。每组 300 条，共 1500 条 Prompt。这样的分层主样本同时覆盖原始基线、现实原型、空间填充、比例表达和边界临界样本，能够在控制调用成本的同时保留主要实验变量。

合并步骤生成 \code{data/experiment_prompts.csv}。为避免不同策略文件中的 \code{prompt_id} 重复，合并表为每条样本增加
\begin{equation}
\mcode{prompt_uid}=\mcode{strategy}\ \Vert\ \mcode{"__"}\ \Vert\ \mcode{prompt_id}.
\end{equation}
其中 \code{prompt_uid} 是 API 采集、断点续跑、特征构建和覆盖率检查的全局 Prompt 标识。

\begin{table}[t]
\centering
\caption{本轮 API 实验 Prompt 来源}
\label{tab:selected-prompts}
\small
\begin{tabularx}{\columnwidth}{p{0.50\columnwidth}rY}
\toprule
来源 & 条数 & 角色 \\
\midrule
\path{baseline} & 300 & 原始基线。\\
\path{real_world_archetype_matrix} & 300 & 现实建筑原型。\\
\path{halton_space_filling} & 300 & 参数空间填充。\\
\path{dimensionless_ratio_design} & 300 & 相对尺度与比例关系。\\
\path{constraint_boundary_solver} & 300 & 临界边界样本。\\
\midrule
合计 & 1500 & 分层主样本。\\
\bottomrule
\end{tabularx}
\end{table}

\subsection{模型集合与重复采样}

实验通过 DF/OpenAI-compatible 网关统一调用模型。模型集合保留最初 4 个模型，并新增 8 个更强或更具代表性的对比模型，共 12 个。每个模型对每条 Prompt 重复采样 3 次，用于区分单次生成波动和模型稳定行为。因此，原始回复规模为
\begin{equation}
1500 \text{ prompts}\times 12 \text{ models}\times 3 \text{ repeats}=54{,}000.
\end{equation}

\begin{table}[t]
\centering
\caption{本轮实验模型集合}
\label{tab:models}
\footnotesize
\begin{tabularx}{\columnwidth}{p{0.46\columnwidth}Y}
\toprule
模型 & 说明 \\
\midrule
\path{gpt-4o} & 原始对比模型。\\
\path{gpt-4o-mini} & 原始轻量对比模型。\\
\path{claude-3-5-sonnet-latest} & 原始 Claude 对比模型。\\
\path{deepseek-chat} & 原始 DeepSeek 对比模型。\\
\path{gpt-5.5} & 新增强模型。\\
\path{claude-opus-4-6} & 新增强模型。\\
\path{claude-sonnet-4-6-thinking} & 新增推理增强模型。\\
\path{gemini-3-flash-preview} & 新增 Gemini 对比模型。\\
\path{grok-4} & 新增 Grok 对比模型。\\
\path{deepseek-v4-pro} & 新增 DeepSeek 强模型。\\
\path{qwen3-max} & 新增 Qwen 强模型。\\
\path{Kimi-K2-Thinking} & 新增 Kimi 推理模型。\\
\bottomrule
\end{tabularx}
\end{table}

\subsection{采集脚本与运行流程}

实验流程由四个脚本串联完成。首先运行 \path{scripts/build_experiment_prompts.py}，从选定 Prompt 文件生成 \path{data/experiment_prompts.csv} 并检查 \code{prompt_uid} 唯一性。随后运行 \path{scripts/collect_experiment_responses.py} 调用模型 API，原始输出写入 \path{data/raw/model_responses_experiment.csv}。接着运行 \path{scripts/build_experiment_features.py}，把原始回复转换为 \path{data/processed/behavior_features_experiment.csv}。最后运行 \path{scripts/analyze_experiment.py}，将策略维度和模型维度统计表写入 \path{outputs_experiment/}。

API 调用使用固定采样参数，当前配置为 \code{temperature=0.2}、\code{max_tokens=900}、\code{timeout_seconds=300}、\code{seed=42}、\code{max_retries=3}。正式运行前先做小样本 smoke test，确认模型可连通、回复非空、最终判断格式可被解析；正式采集时再使用并发 workers 分批推进。

\begin{figure*}[t]
\centering
\begin{tikzpicture}[node distance=0.48cm, every node/.style={font=\scriptsize}]
  \node[dmblock, text width=2.25cm] (src) {Prompt CSVs\\baseline + strategies};
  \node[dmblockgreen, text width=2.25cm, right=of src] (merge) {experiment\\Prompt table\\\code{prompt_uid}};
  \node[dmblockorange, text width=2.05cm, right=of merge] (api) {DF gateway\\12 models};
  \node[dmblock, text width=2.15cm, right=of api] (raw) {raw response\\long table};
  \node[dmblockgreen, text width=2.15cm, right=of raw] (feat) {behavior\\feature table};
  \node[dmblockorange, text width=2.05cm, right=of feat] (mine) {analysis\\tables};
  \draw[dmarrow] (src) -- (merge);
  \draw[dmarrow] (merge) -- (api);
  \draw[dmarrow] (api) -- (raw);
  \draw[dmarrow] (raw) -- (feat);
  \draw[dmarrow] (feat) -- (mine);
  \node[dmnote, below=0.42cm of api, text width=6.0cm] {断点续跑键：$(\mcode{provider},\mcode{model},\mcode{repeat},\mcode{prompt_uid})$};
  \node[dmnote, below=0.42cm of feat, text width=5.0cm] {失败样本保留 \code{status} 与 \code{error_message}，不静默丢弃。};
\end{tikzpicture}
\caption{多模型行为数据采集管线。上游 Prompt 文件首先合并为带有全局唯一键的实验主表，随后经 DF/OpenAI-compatible 网关调用多个模型，最终形成同时保留原始回复、结构化特征和失败记录的长表数据。}
\label{fig:collection-pipeline}
\end{figure*}


\begin{table}[t]
\centering
\caption{实验采集输出文件}
\label{tab:outputs}
\small
\begin{tabularx}{\columnwidth}{p{0.52\columnwidth}Y}
\toprule
路径 & 用途 \\
\midrule
\path{data/experiment_prompts.csv} & 合并后的实验 Prompt 主表。\\
\path{data/raw/model_responses_experiment.csv} & 多模型 API 原始回复表。\\
\path{data/processed/behavior_features_experiment.csv} & 从回复中抽取的行为特征表。\\
\path{outputs_experiment/} & 策略、模型、扰动和层级统计结果。\\
\bottomrule
\end{tabularx}
\end{table}

\subsection{断点续跑、并发与失败记录}

由于完整实验包含 54{,}000 次 API 调用，采集脚本默认按可恢复方式运行。断点续跑键定义为
\begin{equation}
(\mcode{provider},\mcode{model},\mcode{repeat},\mcode{prompt_uid}),
\end{equation}
而不是只使用 \code{prompt_id}。这样即使不同策略中存在相同文件内编号，也不会在续跑时被误判为已完成。

脚本支持 \code{--resume}、\code{--workers}、\code{--log-every} 和 \code{--flush-every}。其中 \code{--workers} 控制并发数量，\code{--log-every} 控制终端进度打印频率，\code{--flush-every} 控制定期写盘间隔。每次调用都会记录 \code{status}、\code{error_message}、\code{latency_seconds} 和完整 \code{response}。失败调用不会被静默丢弃，而是以 \code{status=failed} 进入原始表，便于检查失败是否集中在特定模型、限流或权限问题上。

\begin{table}[t]
\centering
\caption{原始回复表核心字段}
\label{tab:raw-schema}
\small
\begin{tabularx}{\columnwidth}{p{0.42\columnwidth}Y}
\toprule
字段 & 含义 \\
\midrule
\code{strategy}, \code{prompt_uid} & Prompt 来源与全局唯一标识。\\
\code{provider}, \code{model}, \code{repeat} & API 网关、模型名称和重复编号。\\
\code{collected_at} & 回复采集时间。\\
\code{temperature}, \code{max_tokens}, \code{seed} & 调用参数。\\
\code{status}, \code{error_message} & 成功或失败状态及错误信息。\\
\code{latency_seconds} & 单次调用延迟。\\
\code{response} & 模型完整原始回复。\\
\bottomrule
\end{tabularx}
\end{table}
